\chapter{Introduction}%
\label{chp:introduction}

\section{Problem Definition}
Text-to-image generative models have transformed visual content creation, with diffusion-based
architectures achieving high image quality and prompt controllability. Among these, Stable Diffusion \cite{Rombach_2022_CVPR} has emerged as particularly influential due to its open-source availability, widespread deployment, and the extensive ecosystem of tools and fine-tuned variants built upon it.

These systems rely on massive web-scraped datasets to learn the mapping between text and
visual content. Re-LAION-5B, containing over 5 billion image-text pairs harvested from the
internet, represents one of the largest publicly available training corpora for such models
\cite{Schuhmann_2022_Laion5b}. However, the scale and automated nature of web scraping
introduce systematic biases that reflect and potentially amplify societal inequities present
in online content \cite{Birhane_2021_MultiModalDatasetsMisogynyPornography}.

Demographic bias in AI is well-documented across natural language processing
\cite{Bolukbasi_2016_ManComputerProgrammerWoman}, machine learning
\cite{Mehrabi_2022_SurveyBiasFairnessMachine}, and computer vision
\cite{Buolamwini_2018_GenderShades}, yet a gap remains between detecting bias
in generative outputs and understanding the spurious features that cause and
perpetuate it. Current research predominantly measures demographic distributions in generated images against desired benchmarks \cite{Luccioni_2023_AnalyzingSocietalRepresentationsDiffusionModels, Bianchi_2023_TTIGenAmplifiesDemographicStereotypes}, thereby treating models as black boxes and diagnosing symptoms without exposing the underlying mechanisms.

A fundamental yet under explored dimension involves spurious features, which are low-level
non-semantic image characteristics such as colour statistics, texture patterns, edge
distributions, frequency domain properties, and compression artefacts
\cite{Meister_2023_GenderArtifactsinVisualDatasets, Zeng_2024_UnderstandingBiasinLargeScaleVisualDatasets}. When these features correlate with sensitive demographic attributes, models can exploit spurious  statistical regularities rather than semantic content, creating a form of bias that operates beneath standard
evaluation metrics. It also remains unclear how generative training transforms these spurious
features, making it difficult to predict how training data biases manifest in deployed systems.

Existing mitigation approaches typically require expensive interventions: curating balanced
training data \cite{Yang_2020_FairerDatasets}, fairness-constrained training
\cite{Zhang_2018_MitigatingUnwantedBiasesAdversarial}, or fine-tuning on carefully
constructed demographic distributions \cite{Friedrich_2023_FairDiffusionInstructingTextToImage}.
These demand substantial computational resources, creating barriers for researchers working
with pretrained models. Moreover, outcome-focused corrections may balance demographic
distributions without addressing the spurious features that enable biased shortcuts, yielding
superficial fixes that fail to generalise.

This thesis addresses these gaps through a mechanistic investigation of spurious features in
Re-LAION-5B and Stable Diffusion v1.5. Rather than solely measuring outcomes, it applies
established methodologies to expose the spurious features that distinguish dataset sources
and correlate with sensitive attributes, then evaluates targeted inference‐time interventions that address root causes rather than symptoms. Crucially, this investigation reveals that Re‐LAION‐5B and Stable Diffusion v1.5 encode demographic information through qualitatively different mechanisms, a distinction that has direct consequences for how bias is measured, attributed, and mitigated.

\section{Motivation}

\begin{description}
	\item[Influence and Impact:] Re-LAION-5B serves as training data for multiple widely-deployed
	text-to-image models, while Stable Diffusion v1.5 has achieved widespread adoption across
	platforms such as DreamStudio and Hugging Face. Despite this pervasive influence,
	comprehensive analysis of spurious features in Re-LAION-5B and their propagation through
	Stable Diffusion v1.5 remains limited, making the study of bias mechanisms in these systems
	directly relevant to fairness at scale.
	
	\item[Computational Accessibility:] Training or fine-tuning large-scale diffusion models is beyond
	reach for most researchers and practitioners. By focusing on inference-time methods
	applicable to frozen models, this research demonstrates fairness interventions accessible to
	smaller groups and organisations with limited computational budgets.
	
	\item[Reproducibility and Transparency:] The open-source nature of Stable Diffusion v1.5 and the
	research availability of Re-LAION-5B enable findings to be independently verified,
	replicated, and extended, strengthening the credibility and impact of this work.
	
	\item[Beyond Outcome Measurement:] Understanding why biases emerge, and specifically which spurious features enable biased shortcuts, supports interventions that address root causes
	rather than symptoms, enabling solutions that generalise rather than superficially correct
	specific evaluation metrics.
\end{description}

\section{Aims and Objectives}

The overarching aim is to apply established spurious feature detection methodologies to
Re-LAION-5B and Stable Diffusion v1.5, systematically characterise the spurious features
present in each, quantify their correlation with sensitive demographic attributes, and
evaluate the effectiveness of existing inference-time debiasing techniques on this
dataset-model pair. This aim is addressed through the following set of objectives:

\begin{enumerate}
	\item \textbf{Dataset Acquisition, Preparation, and Annotation.}
%	Retrieve images from Re-LAION-5B and generate corresponding images using Stable Diffusion
%	v1.5 across a set of occupation categories. Annotate all images for gender and skin tone
%	using established computer vision models to enable demographic analysis and spurious feature
%	correlation assessment.
	
	Retrieve occupation-conditioned images from Re‐LAION‐5B and generate a corresponding set using Stable Diffusion v1.5, annotating all images for gender and skin tone to support subsequent spurious feature and demographic analysis.
	
	\item \textbf{Spurious Feature Detection.}
%	Implement classifier-based methodologies to detect spurious features distinguishing dataset
%	sources. Apply systematic image transformations to isolate non-semantic features and train
%	discriminative classifiers to characterise dataset-specific spurious patterns.
	
	Detect and characterise spurious features that distinguish Re‐LAION‐5B images from Stable Diffusion v1.5 outputs using transformation-based classifier probing.
	
	\item \textbf{Demographic Correlation Assessment.}
%	Analyse whether identified spurious features exhibit systematic correlations with gender and
%	skin tone in both datasets. Quantify the extent to which spurious features enable demographic
%	classification shortcuts and compare distributions against real-world occupational statistics
%	to identify representation biases.
	
	Probe whether identified spurious features encode demographic information differently across the two sources, and compare occupational demographic distributions against real‐world statistics to assess representation bias and generative amplification.
	
	\item \textbf{Inference-Time Debiasing.}
%	Implement established inference-time debiasing techniques for Stable Diffusion v1.5 without
%	model retraining. Evaluate their effectiveness across demographic representation and spurious
%	feature mitigation, and assess fairness trade-offs to inform practical deployment.

	Implement and evaluate an inference-time debiasing pipeline for Stable Diffusion v1.5, assessing its effectiveness across demographic representation and spurious feature mitigation without model retraining.
	
	\item \textbf{Synthesis of Bias Mechanisms and Mitigation.}
%	Synthesise findings to characterise how spurious features manifest in Re-LAION-5B and
%	propagate through Stable Diffusion v1.5, determining whether the generative process
%	preserves, amplifies, or mitigates training data biases.
	
	Synthesise findings to characterise how spurious features and demographic biases manifest in Re‐LAION‐5B and are transformed by the Stable Diffusion v1.5 generative process.
\end{enumerate}

\section{Document Structure}

The remainder of this dissertation is structured as follows. Chapter~\ref{chp:background_lit} reviews the Re-LAION-5B dataset, Stable Diffusion v1.5, spurious feature learning, demographic annotation frameworks, and inference-time debiasing approaches for text-to-image models. Chapter~\ref{chp:methodology} describes
the image retrieval and generation pipelines, demographic annotation procedure,
transformation-based probing framework, and debiased generation configuration.
Chapter~\ref{chp:evaluation} presents the demographic distribution analysis across
occupations, dataset classification results, demographic attribute separability findings,
and debiasing evaluation. Chapter~\ref{chp:conclusions} summarises the main findings,
discusses their implications for dataset and model assessment, and outlines directions
for future work.